\begin{table}[htbp]
\centering
\caption{Standardized Effect Sizes}
\label{tab:sde}
\begin{adjustbox}{max width=\textwidth}
\begin{tabular}{lcccccc}
\toprule
Outcome & $\hat{\beta}$ & SE & SD($Y$) & SDE & SE(SDE) & Classification \\
\midrule
\multicolumn{7}{l}{\textit{Panel A: Pooled}} \\
SNAP participation rate & -0.0584 & 0.0688 & 7.96 & -0.0073 & 0.0086 & Small negative \\
Poverty rate (placebo) & -0.1384 & 0.0728 & 8.26 & -0.0168 & 0.0088 & Small negative \\
\midrule
\multicolumn{7}{l}{\textit{Panel B: Heterogeneous (sample splits)}} \\
SNAP rate (high-poverty counties) & -0.1045 & 0.1158 & 8.04 & -0.0130 & 0.0144 & Small negative \\
SNAP rate (low-poverty counties) & 0.0104 & 0.0640 & 3.66 & 0.0028 & 0.0175 & Null \\
\bottomrule
\end{tabular}
\end{adjustbox}
\begin{tablenotes}[flushleft]
\small
\item \textit{Notes:} \textbf{Country:} United States. \textbf{Research question:} Does tripling SNAP retailer minimum stocking requirements reduce food stamp participation in counties where small-format retailers dominate the food retail landscape? \textbf{Policy mechanism:} The January 2018 USDA depth-of-stock provision raised the minimum number of staple food items that SNAP-authorized retailers must stock from 12 to 36, creating a binding compliance burden on convenience stores and small groceries that lack shelf space and supply chain capacity to meet the new threshold. \textbf{Outcome definition:} SNAP participation rate, defined as the share of households receiving food stamps/SNAP benefits from ACS table B22003, measured at the county level. \textbf{Treatment:} Continuous; pre-2018 county-level convenience store share of total food retail establishments (standardized, mean zero, unit standard deviation). \textbf{Data:} ACS 5-year estimates (2015--2022 vintages) for 3,163 U.S. counties; County Business Patterns (2015--2016 average) for treatment intensity; 25,304 county-year observations in the balanced panel. \textbf{Method:} Two-way fixed effects (county + state-by-year), standard errors clustered at state level, with randomization inference (500 permutations). \textbf{Sample:} Counties with nonmissing CBP food retail data and population above 100; balanced panel requires all 8 ACS vintages (2015--2022). SDE $= \hat{\beta} / \text{SD}(Y)$ where SD($Y$) is the pre-treatment standard deviation. Classification refers to magnitude, not statistical significance: Large ($|$SDE$| > 0.15$), Moderate ($0.05$--$0.15$), Small ($0.005$--$0.05$), Null ($< 0.005$).
\end{tablenotes}
\end{table}

